\documentclass[11pt]{article}
\usepackage[margin=1in]{geometry}
\usepackage{amsmath, amssymb}

\setlength{\parindent}{0pt}
\setlength{\parskip}{0.6em}

\title{STA4241 Lab 4, Part 1:\\ Two-class LDA implies a logistic model that is linear in $x$}
\author{}
\date{}

\begin{document}
\maketitle

\section*{Claim}

Suppose $Y \in \{0, 1\}$ and the LDA assumptions hold with a single predictor:
\[
X \mid Y = k \;\sim\; N(\mu_k, \sigma^2), \qquad \pi_k = P(Y = k), \qquad k = 0, 1,
\]
where the variance $\sigma^2$ is shared by both classes. Then the log odds of $Y = 1$ are linear in $x$:
\[
\log\!\left( \frac{P(Y = 1 \mid X = x)}{P(Y = 0 \mid X = x)} \right) = \beta_0 + \beta_1 x,
\]
with
\[
\beta_1 = \frac{\mu_1 - \mu_0}{\sigma^2},
\qquad
\beta_0 = \log\!\left(\frac{\pi_1}{\pi_0}\right) - \frac{\mu_1^2 - \mu_0^2}{2\sigma^2}.
\]

\section*{Proof}

\textbf{Step 1: Write the odds using Bayes' rule.}
Let $f_k(x)$ be the density of $X$ in class $k$. Bayes' rule gives
\[
P(Y = 1 \mid X = x) = \frac{\pi_1 f_1(x)}{\pi_1 f_1(x) + \pi_0 f_0(x)},
\qquad
P(Y = 0 \mid X = x) = \frac{\pi_0 f_0(x)}{\pi_1 f_1(x) + \pi_0 f_0(x)}.
\]
The two denominators are identical, so they cancel in the ratio:
\[
\frac{P(Y = 1 \mid X = x)}{P(Y = 0 \mid X = x)} = \frac{\pi_1 f_1(x)}{\pi_0 f_0(x)}.
\]

\textbf{Step 2: Take the log.}
\[
\log\!\left( \frac{P(Y = 1 \mid X = x)}{P(Y = 0 \mid X = x)} \right)
= \log\!\left(\frac{\pi_1}{\pi_0}\right) + \log f_1(x) - \log f_0(x).
\]

\textbf{Step 3: Plug in the normal density.}
For each class,
\[
f_k(x) = \frac{1}{\sqrt{2\pi\sigma^2}} \exp\!\left( -\frac{(x - \mu_k)^2}{2\sigma^2} \right)
\quad\Longrightarrow\quad
\log f_k(x) = -\log\sqrt{2\pi\sigma^2} - \frac{(x - \mu_k)^2}{2\sigma^2}.
\]
Because $\sigma^2$ is shared, the term $-\log\sqrt{2\pi\sigma^2}$ is the same for both classes and cancels in the difference:
\[
\log f_1(x) - \log f_0(x) = \frac{(x - \mu_0)^2 - (x - \mu_1)^2}{2\sigma^2}.
\]

\textbf{Step 4: Expand the squares.}
\begin{align*}
(x - \mu_0)^2 - (x - \mu_1)^2
&= \left( x^2 - 2x\mu_0 + \mu_0^2 \right) - \left( x^2 - 2x\mu_1 + \mu_1^2 \right) \\
&= 2x(\mu_1 - \mu_0) - \left( \mu_1^2 - \mu_0^2 \right).
\end{align*}
The $x^2$ terms cancel. Therefore
\[
\log f_1(x) - \log f_0(x) = \frac{\mu_1 - \mu_0}{\sigma^2}\, x - \frac{\mu_1^2 - \mu_0^2}{2\sigma^2}.
\]

\textbf{Step 5: Collect terms.}
Substituting into Step 2,
\[
\log\!\left( \frac{P(Y = 1 \mid X = x)}{P(Y = 0 \mid X = x)} \right)
= \underbrace{\left[ \log\!\left(\frac{\pi_1}{\pi_0}\right) - \frac{\mu_1^2 - \mu_0^2}{2\sigma^2} \right]}_{\beta_0}
+ \underbrace{\left[ \frac{\mu_1 - \mu_0}{\sigma^2} \right]}_{\beta_1} x ,
\]
which is the logistic regression model $\log\!\big(p(x)/(1 - p(x))\big) = \beta_0 + \beta_1 x$. \hfill $\blacksquare$

\section*{Remarks}

\textbf{1. Same model, different estimation.}
LDA and logistic regression both produce log odds that are linear in $x$. They differ only in how the coefficients are estimated. Logistic regression estimates $(\beta_0, \beta_1)$ directly by maximizing the likelihood of $Y$ given $X$, and assumes nothing about the distribution of $X$. LDA estimates $\mu_0$, $\mu_1$, $\sigma^2$, $\pi_0$, $\pi_1$ and plugs them into the formulas above, which relies on the normality of $X$ within each class.

\textbf{2. Why the shared variance matters.}
The $x^2$ terms cancel in Step 4 only because both classes have the same $\sigma^2$. If the classes had variances $\sigma_1^2 \neq \sigma_0^2$, the $x^2$ term would have coefficient $\frac{1}{2\sigma_0^2} - \frac{1}{2\sigma_1^2} \neq 0$, and the log odds would be quadratic in $x$. That is quadratic discriminant analysis.

\textbf{3. Connection to the discriminant function.}
The LDA discriminant function from lecture is
\[
\delta_k(x) = \frac{x \mu_k}{\sigma^2} - \frac{\mu_k^2}{2\sigma^2} + \log \pi_k .
\]
The log odds derived above equal $\delta_1(x) - \delta_0(x)$. So classifying to class 1 when $\delta_1(x) > \delta_0(x)$ is the same as classifying to class 1 when the log odds are positive, that is, when $P(Y = 1 \mid X = x) > 0.5$.

\end{document}
